\documentclass[letterpaper]{article}
\usepackage[submission]{aaai2026}
\usepackage{times}
\usepackage{helvet}
\usepackage{courier}
\usepackage[hyphens]{url}
\usepackage{graphicx}
\urlstyle{rm}
\def\UrlFont{\rm}
\usepackage{natbib}
\usepackage{caption}
\frenchspacing
\setlength{\pdfpagewidth}{8.5in}
\setlength{\pdfpageheight}{11in}
\pdfinfo{/TemplateVersion (2026.1)}
\usepackage{amsmath,amssymb,booktabs}
% ENGRAM_TAG=v0.2.0
\title{Temporal Perception in AI: Re-Deriving Time for Agent Memory}
\author{Anonymous Submission}
\affiliations{}
\begin{document}
\maketitle
\begin{abstract}
Large language model agents are deployed into a world saturated with time---deadlines, seasons, waiting, change---yet they possess no temporal phenomenology: no felt duration, no experience of the interval between invocations, no native clock. Current practice either ignores this or copies human memory mechanisms (Ebbinghaus decay, recency weighting, sleep-inspired consolidation) whose parameters derive from biological necessities---mortality, metabolism, circadian rhythm---that agents do not have. We argue such borrowing produces \emph{cargo-cult memory}: the form of human mechanisms without their derivational basis. Through phenomenological case analysis we identify systematic disanalogies between human and machine temporality and propose \emph{polytemporal representation}: events located in vectors of heterogeneous time coordinates including wall-clock, monotonic, and subjective $\tau$ (integrated novelty density). Salience at encoding, not age, gates durability; anchor episodes replace calendars as the primary temporal index. We show retention profiles alone induce divergent forgetting (E7 $55/55/35/7$) and real LongMemEval $n{=}50$ gains when decay operates on $\tau$ rather than wall-clock. Our thesis is to borrow wheels, not clocks: reuse cognitive-science mechanisms where the derivation transfers and re-derive the rest from agent-native necessities.
\end{abstract}

% S1 Introduction — Tom vignette 10-case-corpus.md C1 cold winter morning English building, cargo-cult thesis
\section{Introduction}
Tom does not remember the date he first saw his wife---outside the English building at LSU, a cold winter morning. The date is gone; the texture (cold, morning, winter), place, person, and charge are permanent. A sub-second perception outcompetes ten thousand forgotten commutes. This is not a defect of human memory. It is its design: durable retrieval keys are regime-texture, place-context, entities, and valence-magnitude, while wall-clock is a weak reconstructable attribute. Agents inherit none of this structure yet are forced to operate in the same temporal world of deadlines, seasons, and waiting.

Large language model agents possess no temporal phenomenology: no felt duration, no experience of the interval between invocations, no native clock. Current memory practice copies human mechanisms---Ebbinghaus decay $R=e^{-t/S}$, recency weighting, sleep-inspired consolidation---whose parameters derive from biological necessities (mortality, metabolism, circadian rhythm) that agents do not have. As \citet{cheng2026blind} show, agents assume stationary context even when given timestamps, achieving $<65\%$ human alignment on stationarity probes; prompting fails whereas post-training partially fixes. Such behaviour is symptomatic of a deeper problem: we argue it produces \emph{cargo-cult memory}---the form of human mechanisms without their derivational basis.

We borrow wheels, not clocks. Where the cognitive-science derivation transfers (reconstructive recall, spreading activation, complementary learning systems) we reuse it; where it depends on biological contingency we re-derive from agent-native necessities. The case corpus (C1--C8) serves as specification test-vectors: any candidate representation must reproduce the observed encoding-and-recall patterns or fail. From the resulting disanalogy catalog (D1--D10, each linked to falsifiable hypotheses H1--H18) we propose \emph{polytemporal representation}: events located in a heterogeneous \emph{time\_vector} of typed coordinates---wall-clock, monotonic, subjective $\tau$, fuzzy interval, regime phase---with each cognitive mechanism consuming its own projection rather than a privileged axis. Salience at encoding, not age, gates durability; anchor episodes replace flat chronologies; a continuous sensor substrate supplies the experiential stream that a discontinuous inference core lacks. We close with designs runnable on a continuously-observing fleet that pit $\tau$-keyed against wall-clock-keyed memory on decision quality, testing whether re-derived time outperforms borrowed time.

\section{Related Work}
\subsection{Temporal Reasoning}
\citet{fatemi2025tot} (Test of Time, ICLR 2025) and its family---TRAM \citep{tram2024}, TimeBench \citep{chu2024timebench}, TempReason \citep{tan2023tempreason}, TIME \citep{yuan2025time}---decompose temporal competence into semantic and arithmetic components and document systematic failure modes. BORROW: the semantic/arithmetic decomposition and documented failure taxonomy. ADD: all test time-as-content to reason \emph{about}. None test time-as-experience structuring memory. That contrast is our \S2.5 gap claim: nobody re-derives; everybody decorates timestamps. A second strand studies whether models know time passes at all: token-time versus wall-clock-time as distinct measurement domains \citep{discrete2025continuous}, with evidence models use token-length as a duration proxy and adapt under time pressure, and surveys of temporal sense-making in LLMs \citep{cacm2026time}. BORROW: token-time vs.\ wall-clock as measurement domains. ADD: token-time remains chronometry; our $\tau$ is experiential density, and continuity/discontinuity (C8) is absent there entirely. \citet{cheng2026blind} provide the empirical ground for our cargo-cult critique: agents are temporally blind even with timestamps. BORROW: proof that prompting fails and post-training partially fixes. ADD: they align agents \emph{to} human time as a tool-use skill; we argue the deeper fix is a native time representation with alignment as translation between coordinate systems. \citet{forever2026} is our closest ally: model-time $\neq$ wall-clock, with replay scheduled by accumulated optimizer-update magnitude. BORROW: native-time principle. ADD: they schedule continual-learning replay by gradients; we extend native-time to episodic memory and define $\tau$ from experience density (novelty, prediction-error) available to frozen models.

\subsection{Agent Memory}
\citet{packer2023memgpt} cast LLM memory as an operating system with self-directed paging. BORROW: OS-paging analogy and self-directed archival. ADD: no consolidation, no decay policy, salience as mid-task attention; hoarding tendency. \citet{zhong2023memorybank} adopt $R=e^{-t/S}$ as the wall-clock baseline (E1). We retain the functional form but critique the variable: $t$ is wall-clock---a cargo-cult parameter whose form derives from biological consolidation dynamics the agent lacks. BORROW: $R=e^{-t/S}$ wall-clock baseline our E1 must beat. ADD: $\tau$ substitutes for $t$ and is the experimental contrast. \citet{park2023generative} introduce recency $\times$ importance $\times$ relevance retrieval. BORROW: triad structure. ADD: importance is a scored attribute, not an encoding-time durability tier; no anchors, no constellations, no reconstruction grading. \citet{gutierrez2024hipporag} implement hippocampal indexing via personalized PageRank concept graphs. BORROW: HT-indexing theory. ADD: retrieval-only---borrows the index, skips the lifecycle (encoding gates, consolidation, forgetting). \citet{liu2024actr} bring ACT-R base-level activation decay and noise to agent memory. BORROW: activation dynamics. Cognitive substrates such as Engram \citep{engram2024} compose ACT-R activation, Ebbinghaus decay, Hebbian links, and interoceptive hubs (proof the stack composes); BORROW: substrate composition. ADD: drives are hardcoded constants; ours are derived from the disanalogy catalog (what necessity substitutes for metabolism?). Surveys \citep{niu2024survey} and difference literature on embodiment confirm LLMs differ because they are not embodied, framing our contribution as re-derivation rather than decoration.

\subsection{Cognitive Foundations}
We borrow wheels from: James (stream of consciousness, \S7 continuity) \citep{james1890principles}; Ebbinghaus (decay curve origin) \citep{ebbinghaus1885memory}; Bartlett (reconstructive recall; schema-consistent intrusions, H5) \citep{bartlett1932remembering}; Atkinson \& Shiffrin and Tulving (store taxonomy, episodic/semantic split) \citep{atkinson1968human,tulving1972episodic,tulving1983elements}; Collins \& Loftus (spreading activation, H10) \citep{collins1975spreading}; Gibbon (Scalar Expectancy Theory, scalar variability $\to$ fuzzy intervals H9) \citep{gibbon1977scalar}; McClelland et al.\ (Complementary Learning Systems, fast/slow learners $\to$ curator/consolidator) \citep{mcclelland1995cls}; Conway \& Pleydell-Pearce (Self-Memory System, lifetime periods $\supset$ general events $\supset$ specific knowledge $\to$ H7) \citep{conway2000sms}; Zacks \& Swallow (event segmentation at prediction-error boundaries $\to$ schema fusion) \citep{zacks2007event}; Bergson (lived vs.\ spatialized time) \citep{bergson1889duree}; Schacter (adaptive forgetting) \citep{schacter2001seven}; Locke (psychological continuity) \citep{locke1690essay}; Schmidhuber (compression-progress, curiosity/boredom $\to$ $\tau$ and D8) \citep{schmidhuber1991curiosity}. Curiosity and boredom formalized as compression progress give $\tau$ its definition and justify boredom as a scheduling signal rather than a defect.

\section{Disanalogy Catalog}
Each disanalogy states a human mechanism, its biological derivation, the agent substitute, architectural consequence, and testable prediction (H*). Together they source C1--C8 into falsifiable claims.

\paragraph{D1 Texture versus timestamp (C1).} Humans encode regime-texture (cold, morning, season, place) durably; wall-clock is transient and reconstructed. Biology: episodic binding prioritizes context useful for survival. Substitute: index by typed regime coordinates (cycle, phase, day-class) plus wall-clock as payload attribute. Consequence: retrieval by texture beats date-range on human-style probes. $\to$ H1 \citep{cheng2026blind}.

\paragraph{D2 Salience-gated durability (C1--C2).} Durability tier is set at write-time by surprise $\times$ goal-weight; flashbulb memories become exempt from decay. Biology: amygdala-hippocampal consolidation gated by arousal. Substitute: salience $s\in[0,1]$ at encoding selects tier (normal/flashbulb/protected) and initial strength $S$; low-salience never enters the episode store. $\to$ H2.

\paragraph{D3 Anchored constellations versus flat chronologies (C5).} Humans store landmark-relative edges (before/during/after an anchor) not absolute dates; the calendar is reconstructed by graph traversal. Biology: hierarchical autobiographical memory with anchor hubs. Substitute: constellation graph over anchor episodes with typed edges. $\to$ H3.

\paragraph{D4 Schema fusion versus episodic hoarding (C3).} Repeated near-identical episodes compress into a schema-trace; only prediction-errors retain episode status. Consequence: storage grows $\sim\log$ with experience; recall of ``what usually happens'' improves while per-instance recall of routine correctly fails. $\to$ H4, H14.

\paragraph{D5 Reconstruction versus lookup (C2).} Recall is generative inference from cues plus schemas, graded by confidence---not record lookup. Errors cluster as schema-consistent intrusions; confidence tracks accuracy. $\to$ H5, H6 \citep{bartlett1932remembering}.

\paragraph{D6 Subjective time density versus uniform seconds (C6, C8).} Felt duration dilates with novelty, prediction-error, and engagement; decay should operate on $\tau$ not wall-clock. $d\tau/dt = f(\text{novelty-rate}, \text{prediction-error}, \text{engagement})$. $\tau$-keyed memory beats wall-clock-keyed memory on decision-quality probes across quiet versus eventful stretches of equal wall duration. $\to$ H-$\tau$ (E1).

\paragraph{D7 Continuity versus discontinuity (C8).} Human consciousness is an uninterrupted sensorimotor stream; agent inference is discrete and gapped. Biology: continuity survives sleep as a mode of the stream, not absence. Substitute: couple a discontinuous inference core to a continuous sensor substrate (bot fleet) and mark witnessed versus reconstructed change. $\to$ H11--H13 \citep{james1890principles,bergson1889duree}.

\paragraph{D8 Drives versus drives-absent (boredom/fatigue as signals).} Human boredom and fatigue schedule exploration and rest. Agents lack metabolism but need analogous scheduling signals for encoding and consolidation cadence. Substitute: curiosity/boredom as compression-progress signals driving gate thresholds and $\tau$ accumulation. $\to$ H8/H-$\tau$ \citep{schmidhuber1991curiosity}.

\paragraph{D9 Landmark-relative dating versus absolute dates (C5, C7).} Even when absolute time is stored, humans answer ``when'' by traversing anchors (``right after the flood''). Substitute: qualifier-typed anchor edges plus fuzzy intervals whose width encodes confidence. $\to$ H3, H9 \citep{gibbon1977scalar}.

\paragraph{D10 Precision as confidence (C7).} Humans store intervals $[t_{\text{lo}},t_{\text{hi}}]$ whose width equals encoding confidence; recall returns answers at stored resolution, never manufactured precision. The wallet lost ``between noon and sunset'' correctly reports an interval. $\to$ H9. Calibrated-confidence scoring predicts that reported interval widths track actual temporal error.

\section{Polytemporal Representation}
Every episode carries a heterogeneous \emph{time\_vector} of typed, indexed coordinates (Table~\ref{tab:polytemporal}). No single axis is privileged; each mechanism consumes its own projection. The polytemporal schema enforces this by giving coordinates real types and indexes---no JSONB on any query path ($\S$25).

\begin{table}[t]\centering\begin{tabular}{ll}\toprule Coordinate & Type \\ \midrule wall & timestamptz \\ tau & double \\ fuzz & tstzrange \\ \bottomrule\end{tabular}\caption{Polytemporal vector (typed, indexed).}\label{tab:polytemporal}\end{table}

Coordinates per \texttt{episode} (25-polytemporal-schema.md): \texttt{wall} \texttt{timestamptz} (human interface; payload attribute), \texttt{mono} \texttt{bigint} (capture-stream ordering), \texttt{tau} \texttt{double precision} NOT NULL (subjective time $\int$ novelty-density), \texttt{fuzz} \texttt{tstzrange} $[t_{\text{lo}},t_{\text{hi}}]$ with width $=$ encoding confidence, \texttt{cycle\_id} + \texttt{phase\_idx} (regime position), \texttt{day\_class} and \texttt{precision\_src} as lookup-typed smallints, plus \texttt{salience}, \texttt{tier}, \texttt{strength} (decay $S$; flashbulb $\Rightarrow$ \texttt{Infinity}), and \texttt{provenance} (witnessed/boundary/reconstructed). Indexes: \texttt{episode\_tau\_idx} on \texttt{tau}, \texttt{episode\_fuzz\_idx} GiST on \texttt{fuzz}, \texttt{episode\_regime\_idx} on (\texttt{cycle\_id}, \texttt{phase\_idx}).

Subjective time is defined as $\tau(t) = \int_{0}^{t} f(\text{novelty-rate}(u), \text{prediction-error}(u), \text{engagement}(u))\,du$, i.e.\ $d\tau/dt = f(\cdot)$, so that $d\tau/dt\approx 0$ on uneventful stretches and spikes on surprise \citep{schmidhuber1991curiosity}. All decay, reinforcement, and consolidation cadence operate on $\tau$; the canonical E1 recall query ranks by $\text{strength}\cdot \exp(-(\tau_{\text{now}}-\tau)/S)$. Interval probes use containment $\texttt{fuzz} \&\& \$window$ (H9); constellation walks use recursive spread over \texttt{anchor\_edge} (H3); regime slicing filters on (\texttt{cycle\_id}, \texttt{phase\_idx}). Fuzz intervals implement H9: positions are labeled intervals whose reported width materializes the scalar variability of temporal encoding \citep{gibbon1977scalar}, and recall is graded by whether answers respect that resolution. Gate-at-encoding (H14) ensures $d\tau/dt$ advances only on surprise; repeated near-identical episodes never enter the store individually but reinforce the schema-trace (H4/H15 resolution ladder), while escalation bursts (H16) temporarily hold the gate open on high-arousal triggers and deferred significance promotes tiers from downstream reference counts.

\bibliography{refs}
\end{document}
